\documentclass[10pt,a4paper]{article}
\usepackage{amsmath,amssymb,amsthm,graphicx}
\title{Logspace-uniform logarithmic-depth approximation of substochastic matrix powers and BPL inclusion}
\author{Kuan Cheng, Yichuan Wang}
\date{doi:10.4230/LIPIcs.CCC.2024.32; CC BY4.0}
\begin{document}
\maketitle
\noindent\textit{Allowed standard prerequisites (editorial).}\begin{itemize}
\item Standard induced L1 matrix norm inequalities and telescoping difference-of-powers bound. Author Theorem5, exact statements retained.
\item Standard probabilistic-space machine transition-matrix representation. Author introduction p.32:2, exact paragraph retained; author Corollary21 is the immediate stated consequence.
\item Ajtai constant-depth approximate majority and convolution pairwise hashing. Author Lemma15 cites[2], Definition12/Theorem13 cites[22]; statements and exact construction retained with complete local approximate-counting proof.
\end{itemize}\medskip
\section*{Author circuit and hashing definitions}
\paragraph{Definition 8 (AC circuit).} AC circuit is a circuit with input gates, NOT gates, unbounded fan-in AND/OR gates, and (possibly more than one) output gates. The size of a circuit is defined by the number of AND/OR gates. The depth of a circuit is defined by the largest number of AND/OR gates on any path from an input gate to an output gate.
\paragraph{Definition 9 (L-uniform AC circuit family).} For functions $S,d:\mathbb Z^+\to\mathbb R^+$, we say a collection of circuits $\{C_n\}_{n\in\mathbb Z^+}$ is an $S$-size $d$-depth L-uniform AC circuit family, if each $C_n$ has size $\le S(n)$ and depth $\le d(n)$, and given the binary representation of $n$, the description of $C_n$ can be computed in uniform $O(\log n)$-space.
We need to mention that the number of input gates in $C_n$ is not necessarily $n$. Also note that since we can encode a tuple of $O(1)$ many integers to a single integer, we can also consider circuit collections with a tuple of integers as an index.
\paragraph{Definition 10 (Complexity Class L-AC$^1$).} We say a language $L$ is in the class L-AC$^1$ if there exists a $\operatorname{poly}(n)$-size $O(\log n)$-depth L-uniform AC circuit family $\{C_n\}$ such that $C_n$ computes $L$ on $n$-bit inputs.
\paragraph{Definition 11 (GapMaj).} For $n\in\mathbb Z^+$ and $a,b\in\mathbb R$ such that $0\le a<b\le n$, define the promise problem $\operatorname{GapMaj}[a,b]$ on $n$ bits as follow:
\[\operatorname{GapMaj}[a,b](x_1,\ldots,x_n):=\begin{cases}\mathrm{YES}&\text{if }x_1,\ldots,x_n\text{ contains }\ge b\text{ 1's}\\\mathrm{NO}&\text{if }x_1,\ldots,x_n\text{ contains }\le a\text{ 1's}\\\bot&\text{otherwise.}\end{cases}\]
We will use pairwise independent hash function as a tool for approximate counting in AC. We shall use the following construction based on convolution, which was also used in [22].
\paragraph{Definition 12 (Convolution-Based Pairwise Independent Hash Function).} Suppose $m$ is a power of 2. Define $H_m:[m^3]\times[m]\to[m]$ by: for $(k,x)\in[m^3]\times[m]$, let $x_1\cdots x_{\log m}$ be the binary representation of $x-1$, let $a_1\cdots a_{2\log m}b_1\cdots b_{\log m}$ be the binary representation of $k-1$, let $y_j:=(\sum_{i=1}^{\log m}a_{i+j}x_i+b_j)\bmod2$ for $j\in[\log m]$, then define $H_m(k,x)$ by letting $y_1\cdots y_{\log m}$ be the binary representation of $H_m(k,x)-1$.
\paragraph{Theorem 13.} $H_m$ is Pairwise Independent Hash Function in the following sense: for any $1\le i<j\le m$, when $k$ is sampled from the uniform distribution over $[m^3]$, the joint distribution of $(H_m(k,i),H_m(k,j))$ is identical to the uniform distribution over $[m]\times[m]$.

\section*{Author matrix conventions and complexity context}
One may view derandomizing BPL as the problem of approximating powers of substochastic matrices. For a TM with $s$ bits of memory, one can label all its states by elements in $[2^s]$. We can define $A\in\mathbb R^{2^s\times2^s}$ to be its transition matrix in the sense that $A_{i,j}$ is the probability that the machine moves from state $i$ to state $j$ in one step. Note that it must arrive at an accept or reject state in $2^s$ steps, so we only need to approximate $A^{2^s}$.
\paragraph{Definition 4 ($L_1$-norm).} Define the $L_1$-norm of a vector $(x_1,\ldots,x_n)^\top\in\mathbb R^n$ to be
\[\|(x_1,\ldots,x_n)^\top\|_1:=|x_1|+\cdots+|x_n|.\]
Define the $L_1$-norm of a matrix $A\in\mathbb R^{n\times n}$ to be
\[\|A\|_1:=\sup_{x\in\mathbb R^n}\frac{\|Ax\|_1}{\|x\|_1}=\max_{1\le j\le n}\{|A_{1,j}|+|A_{2,j}|+\cdots+|A_{n,j}|\}.\]
\paragraph{Theorem 5.} For any $A,B\in\mathbb R^{n\times n}$, we have:
\begin{enumerate}\item $\|A+B\|_1\le\|A\|_1+\|B\|_1$;\item $\|AB\|_1\le\|A\|_1\|B\|_1$;\item If $\|A\|_1,\|B\|_1\le1$, then for any $p\in\mathbb Z^+$, $\|A^p-B^p\|_1\le p\|A-B\|_1$.\end{enumerate}
\paragraph{Definition 6 (Non-negative Matrix).} We say a matrix is non-negative if each of its entry is non-negative.
\paragraph{Definition 7 (Substochastic Matrix).} We say a matrix $A\in\mathbb R^{n\times n}$ is a substochastic matrix if $A$ is non-negative and $\|A\|_1\le1$.
For simplicity, we always assume that the size of a substochastic matrix is a power of 2. To represent a substochastic matrix, we independently represent each entry in binary, accurate to $100\log n$ decimal places.

\paragraph{Theorem 20.} Let $n$ be a power of 2. Then there exists a $\operatorname{poly}(n)$-size $O(\log n)$-depth L-uniform AC circuit family $\{\mathcal M_n\}$\footnote{We require that given $n$, description of $\mathcal M_n$ can be computed in space $O(\log n)$.} such that on input a substochastic matrix $A\in\mathbb R^{n\times n}$, $\mathcal M_n$ outputs a substochastic matrix $M\in\mathbb R^{n\times n}$ such that $\|M-A^n\|_1\le1/n$.
\paragraph{Corollary 21.} $\mathrm{BPL}\subseteq\mathrm{L}\text{-}\mathrm{AC}^1$.

\paragraph{Theorem 14.} Let $n,a,b\in\mathbb N$ such that $0\le a<b\le n$. Then there exists a $\operatorname{poly}(n)$-size $O\left(\frac{\log\frac b{b-a}}{\log\log n}+1\right)$-depth L-uniform AC circuit family $\{C_{n,a,b}\}$ that computes $\operatorname{GapMaj}[a,b]$ on $n$ bits.

The proof depends on the next few Lemmas.
\paragraph{Lemma 15 ([2]).} Let $n\in\mathbb Z^+$. Then there exists $\operatorname{poly}(n)$-size $O(1)$-depth L-uniform AC circuit family $\{C_n^{(0)}\}$ that computes $\operatorname{GapMaj}[n/3,2n/3]$ on $n$ bits.
\paragraph{Lemma 16 (Exact Counting).} Let $n,\ell\in\mathbb Z^+$ such that $n\ge\ell$. Then there exists a $\operatorname{poly}(n)$-size $O\left(\frac{\log\ell}{\log\log n}+1\right)$-depth L-uniform AC circuit family $\{E_{n,\ell}\}$ such that on $\ell$ bits of input, $E_{n,\ell}$ outputs the exact number of 1's over the input bits, in binary form.
We remark that in Lemma 16, $n$ is only used to bound the size of the circuit. Also, $n,\ell$ are not necessarily polynomially related.
\begin{proof}
We only need to show how to compute sum of $O(\sqrt{\log n})$ many $O(\log n)$-bit\footnote{Here ``$O(\log n)$-bit integers'' refers to integers which has $O(\log n)$-bits in its binary representation.} non-negative integers in $O(1)$-depth, then by divide-and-conquer we can compute sum of $\ell$ bits in $O\left(\frac{\log\ell}{\log\log n}+1\right)$-depth.
View the $O(\log n)$-bit integers as $2^{\lfloor\sqrt{\log n}\rfloor}$-base $O(\sqrt{\log n})$-digit integers. We use the grade-school algorithm to sum $O(\sqrt{\log n})$ integers as follow. We first guess the result and all carry-bits, which involve at most $O(\sqrt{\log n})\cdot O\left(\log\left(\sqrt{\log n}\cdot2^{\lfloor\sqrt{\log n}\rfloor}\right)\right)=O(\log n)$ bits, and thus has at most $\operatorname{poly}(n)$ choices. Then we can apply a local check on each digit, each local check involves at most $O(\log n)$ bits, and thus deciding whether all local checks are passed can be computed in $O(1)$-depth. Then we can take the result of the only guess that passes all local checks. The total cost is $O(1)$-depth.
\end{proof}
\paragraph{Lemma 17.} Let $n,a,b\in\mathbb N$ such that $0\le a<b\le n$. Then there exists a $\operatorname{poly}(n)$-size $O\left(\frac{\log\frac n{b-a}}{\log\log n}+1\right)$-depth L-uniform AC circuit family $\{C_{n,a,b}^{(1)}\}$ that computes $\operatorname{GapMaj}[a,b]$ on $n$ bits.
\begin{proof}
Only consider the case that $n$ is a power of 2, otherwise we can use a simple padding argument. By Lemma 15, it suffices to show how to reduce $\operatorname{GapMaj}[a,b]$ on $n$ bits to $\operatorname{GapMaj}[n^3/3,2n^3/3]$ on $n^3$ bits, via a $\operatorname{poly}(n)$-size $O\left(\frac{\log\frac n{b-a}}{\log\log n}+1\right)$-depth L-uniform AC circuit.
If $b-a\le4\sqrt n$ then we can directly compute the number of 1's exactly via Lemma 16. Below we only consider $b-a>4\sqrt n$.
Let $\ell:=\left\lceil\frac{12n^2}{(b-a)^2}\right\rceil$. Suppose the $\operatorname{GapMaj}[a,b]$ instance is $x_1,x_2,\ldots,x_n$. Let $H_n$ be the hash function defined in Definition 12. Define $y_1,\ldots,y_{n^3}$ as follow: for $i\in[n^3]$, let $y_i$ be 1 if at least $\frac{a+b}{2n}$ fraction of $x_{H_n(i,1)},\ldots,x_{H_n(i,\ell)}$ is 1, otherwise let $y_i$ be 0. Note that $y_1,\ldots,y_{n^3}$ can be computed via a $\operatorname{poly}(n)$-size $O\left(\frac{\log\ell}{\log\log n}+1\right)$-depth L-uniform AC circuit, by Lemma 16. Here $O\left(\frac{\log\ell}{\log\log n}+1\right)=O\left(\frac{\log\frac n{b-a}}{\log\log n}+1\right)$.
Let's do some simple calculations. Assume $p$ fraction of $x_1,\ldots,x_n$ is 1. Let $S_i$ be number of 1's in $x_{H_n(i,1)},\ldots,x_{H_n(i,\ell)}$. Then we have $\mathbb E_{i\sim[n^3]}[S_i]=p\ell$ and $\operatorname{Var}_{i\sim[n^3]}[S_i]\le\ell$. So if $p\le\frac an$, then
\[\Pr_{i\sim[n^3]}\left[S_i\ge\ell\cdot\frac{a+b}{2n}\right]\le\frac\ell{(\ell\cdot\frac{b-a}{2n})^2}=\frac{4n^2}{\ell(b-a)^2}\le\frac13.\]
Similarly if $p\ge\frac bn$ then
\[\Pr_{i\sim[n^3]}\left[S_i\le\ell\cdot\frac{a+b}{2n}\right]\le\frac13.\]
This means if $x_1,\ldots,x_n$ is YES/NO instance of $\operatorname{GapMaj}[a,b]$, then $y_1,\ldots,y_{n^3}$ is YES/NO instance of $\operatorname{GapMaj}[n^3/3,2n^3/3]$. The reduction is completed.
\end{proof}
\begin{proof}[Proof of Theorem 14]
We will try to reduce to Lemma 17. Suppose the $\operatorname{GapMaj}[a,b]$ instance is $x_1,x_2,\ldots,x_n$. We only consider the case $n$ is a power of 2, otherwise use a simple padding argument. We only consider the case $10\left(\frac b{b-a}\right)^2<\frac n{b-a}$ (or equivalently, $n(b-a)>10b^2$), otherwise we can directly apply Lemma 17.
Let $\ell:=\left\lfloor\frac{n(b-a)}{2b^2}\right\rfloor$. For $i\in[n^3]$, let $y_i:=x_{H_n(i,1)}\vee\cdots\vee x_{H_n(i,\ell)}$, here $H_n$ is the hash function defined in Definition 12. Then $y_1,\ldots,y_{n^3}$ can be computed via $\operatorname{poly}(n)$-size $O(1)$-depth L-uniform AC circuit.
Assume $p$ fraction of $x_1,\ldots,x_n$ is 1. Let $S_i$ be number of 1's in $x_{H_n(i,1)},\ldots,x_{H_n(i,\ell)}$. Then we have $\mathbb E_{i\sim[n^3]}[S_i]=p\ell$ and $\mathbb E_{i\sim[n^3]}[S_i^2]=\ell(\ell-1)p^2+\ell p\le\ell p+\ell^2 p^2$. Thus by
\[\frac{\mathbb E_{i\sim[n^3]}[S_i]^2}{\mathbb E_{i\sim[n^3]}[S_i^2]}\le\Pr_{i\sim[n^3]}[S_i\ge1]\le\mathbb E_{i\sim[n^3]}[S_i]\]
we know: if $p\le\frac an$, then $\Pr_{i\sim[n^3]}[S_i\ge1]\le\frac{\ell a}n$; if $p\ge\frac bn$, then
\[\Pr_{i\sim[n^3]}[S_i\ge1]\ge\frac{(\frac{\ell b}n)^2}{\frac{\ell b}n+(\frac{\ell b}n)^2}\ge\frac{\ell b}n-\left(\frac{\ell b}n\right)^2.\]
To summarize, if $x_1,\ldots,x_n$ is YES/NO instance of $\operatorname{GapMaj}[a,b]$, then $y_1,\ldots,y_{n^3}$ is YES/NO instance of
\[\operatorname{GapMaj}\left[\left\lfloor n^3\cdot\frac{\ell a}n\right\rfloor,\left\lceil n^3\cdot\left(\frac{\ell b}n-\left(\frac{\ell b}n\right)^2\right)\right\rceil\right].\]
Finally we observe that
\[\left(\frac{\ell b}n-\left(\frac{\ell b}n\right)^2\right)-\frac{\ell a}n=\ell\cdot\left(\frac{b-a}n-\frac{\ell b^2}{n^2}\right)\ge\frac{n(b-a)}{3b^2}\cdot\frac{b-a}{2n}=\frac{(b-a)^2}{6b^2}.\]
Thus by Lemma 17, $\operatorname{GapMaj}\left[\left\lfloor n^3\cdot\frac{\ell a}n\right\rfloor,\left\lceil n^3\cdot\left(\frac{\ell b}n-\left(\frac{\ell b}n\right)^2\right)\right\rceil\right]$ over $n^3$ bits can be computed via a $\operatorname{poly}(n)$-size $O\left(\frac{\log\frac b{b-a}}{\log\log n}+1\right)$-depth L-uniform AC circuit.
\end{proof}

\section*{4 The Iteration Method}
In this section, we will introduce the iteration step, which is the core of our proof of $\mathrm{BPL}\subseteq\mathrm{L}\text{-}\mathrm{AC}^1$.
\paragraph{Theorem 18 (The Iteration).} Let $A\in\mathbb R^{n\times n}$ be a substochastic matrix and $k,t\in\mathbb Z^+$ such that $\log n\ge k\ge t$. Suppose substochastic matrices $B_0,\ldots,B_{k-1}$ are approximations of $A^{2^0},\ldots,A^{2^{k-1}}$ such that $\|B_i-A^{2^i}\|_1\le\varepsilon_i$ for $i=1,2,\ldots,k-1$. Define
\begin{align*}
C:=&-\sum_{i=1}^{t-1}\sum_{\substack{\{j_1<\cdots<j_p\}\uplus\{j_1'<\cdots<j_q'\}\\=\{k-1,k-2,\ldots,k-i+1\}}}B_{j_p}\cdots B_{j_1}B_{k-i}^2 B_{j_1'}\cdots B_{j_q'}\\
&+\sum_{\substack{\{j_1<\cdots<j_p\}\uplus\{j_1'<\cdots<j_q'\}\\=\{k-1,k-2,\ldots,k-t+1\}}}B_{j_p}\cdots B_{j_1}B_{k-t}^2 B_{j_1'}\cdots B_{j_q'}.
\end{align*}
Then
\[\|C-A^{2^k}\|_1\le\sum_{i=1}^{t-1}2^{i-1}\varepsilon_{k-i}^2+2^t\varepsilon_{k-t}.\]
\begin{proof}
Note that
\begin{align*}
C-A^{2^k}=&-\sum_{i=1}^{t-1}\sum_{\substack{\{j_1<\cdots<j_p\}\uplus\{j_1'<\cdots<j_q'\}\\=\{k-1,k-2,\ldots,k-i+1\}}}B_{j_p}\cdots B_{j_1}(A^{2^{k-i}}-B_{k-i})^2B_{j_1'}\cdots B_{j_q'}\\
&-\sum_{\substack{\{j_1<\cdots<j_p\}\uplus\{j_1'<\cdots<j_q'\}\\=\{k-1,k-2,\ldots,k-t+1\}}}B_{j_p}\cdots B_{j_1}(A^{2^{k-t+1}}-B_{k-t}^2)B_{j_1'}\cdots B_{j_q'}.
\end{align*}
So by Theorem 5,
\begin{align*}\|C-A^{2^k}\|_1&\le\sum_{i=1}^{t-1}2^{i-1}\|A^{2^{k-i}}-B_{k-i}\|_1^2+2^t\|A^{2^{k-t}}-B_{k-t}\|_1\\&\le\sum_{i=1}^{t-1}2^{i-1}\varepsilon_{k-i}^2+2^t\varepsilon_{k-t}.\end{align*}
\end{proof}
\paragraph{Theorem 19 (Computing the Iteration).} Let $n,k,t,A,B_0,\ldots,B_{k-1},\varepsilon_0,\ldots,\varepsilon_{k-1},C$ be as defined in Theorem 18. Let $d$ be an integer such that $4\log n\ge d\ge t/10$. Then there exists a $\operatorname{poly}(n)$-size $O(d)$-depth L-uniform AC circuit family $\{\mathcal I_{n,k,t,d}\}$ such that on inputs $B_{k-t},\ldots,B_{k-1}$, if
\[\sum_{i=1}^{t-1}2^{i-1}\varepsilon_{k-i}^2+2^t\varepsilon_{k-t}\le\frac1{2^{d+2}}\]
is satisfied, then $\mathcal I_{n,k,t,d}$ outputs a substochastic matrix $C'$ such that $\|C'-A^{2^k}\|_1\le1/2^d$.
The intuition behind Theorem 19 is that to approximately compute $C$, all arithmetic operations only need a multiplicative accuracy of $1/2^{\Theta(d)}$. This can be done efficiently by L-uniform AC circuit by Theorem 14.
\begin{proof}[Proof of Theorem 19]
We observe that $C$ is the sum of $2^{t-1}$ ``+'' terms and $2^{t-1}-1$ ``-'' terms, and each term is a multiplication of not more than $t+1$ substochastic matrices. We will first show how to approximate the multiplication of substochastic matrices and then show how to approximate their sum.
To approximate $Z:=XY$ for two substochastic matrices $X,Y$, we only need to approximate $\sum_{r=1}^n X_{i,r}Y_{r,j}$ for each pair $(i,j)\in[n]^2$. We first represent each entry $X_{i,r},Y_{r,j}$ using $n^{100}$ bits such that fraction of 1's in these $n^{100}$ bits is equal to the entry, then use a layer of AND gate to represent each $X_{i,r}Y_{r,j}$ using fraction of 1's in $n^{200}$ bits, and then represent each $\frac1n\sum_{r=1}^n X_{i,r}Y_{r,j}$ using fraction of 1's in $n^{201}$ bits. Then we invoke $C_{n^{201},\ell,\lceil\ell(1+1/2^{20d+10})\rceil}$ (as defined in Theorem 14, which has depth $\le O(\frac d{\log\log n}+1)\le O(\frac d{\log(t+1)})$)\footnote{In Theorem 14 we take $(a,b)=(\ell,\lceil\ell(1+1/2^{20d+10})\rceil)$, and then $\log\frac b{b-a}\le O(d)$.} for $\ell=1,2,\ldots,n^{200}$ over these $n^{201}$ bits. Suppose $\ell_0$ is the smallest index such that $C_{n^{201},\ell_0,\lceil\ell_0(1+1/2^{20d+10})\rceil}$ outputs 0, then we have
\[\frac{\ell_0-1}{n^{200}}<Z_{i,j}<\frac{\ell_0(1+\frac1{2^{20d+10}})}{n^{200}}\]
and thus\footnote{Since $n^{200}Z_{i,j}$ is an integer, we have $\frac{\ell_0-1}{n^{200}}<Z_{i,j}\Longrightarrow\frac{\ell_0}{n^{200}}\le Z_{i,j}$.}
\[\frac{Z_{i,j}}{1+\frac1{2^{20d+10}}}-\frac1{n^{100}}\le\frac1{n^{100}}\left[\frac{\ell_0}{n^{100}}\right]\le Z_{i,j}.\]
Use $[\ell_0/n^{100}]/n^{100}$ as an approximation of $Z_{i,j}$, then we obtain an approximation $\widetilde Z$ of $Z$ such that $Z-\widetilde Z$ is non-negative and $\widetilde Z$ is substochastic and $\|Z-\widetilde Z\|_1\le1/2^{20d+10}+1/n^{99}$. We need to be careful that here we need a multiplicative small error on each entry and thus we need to strengthen Lemma 17 to Theorem 14.
Then multiplication of not more than $t+1$ substochastic matrices can be computed via $O(\log(t+1))$ layers of multiplication of two matrices. Recall that multiplying two matrices uses $O(\frac d{\log(t+1)})$-depth and has additive error $1/2^{20d+10}+1/n^{99}$. So the total depth for computing multiplication of not more than $t+1$ substochastic matrices is $O(d)$ and the total error is $\le t(1/2^{20d+10}+1/n^{99})\le1/2^{19d+5}$.
To summarize, suppose $C=-\sum_{i=1}^{2^{t-1}-1}D_i+\sum_{i=1}^{2^{t-1}}D_i'$, here each $D_i,D_i'$ is multiplication of some substochastic matrices. Then we can compute their approximations $\widetilde D_i,\widetilde D_i'$ in $O(d)$ depth such that $\|D_i-\widetilde D_i\|_1\le1/2^{19d+5}$ and $\|D_i'-\widetilde D_i'\|_1\le1/2^{19d+5}$.
We approximate $\frac1{2^{t-1}}\sum_{i=1}^{2^{t-1}-1}\widetilde D_i$ and $\frac1{2^{t-1}}\sum_{i=1}^{2^{t-1}}\widetilde D_i'$. Use the similar idea as summing $\frac1n\sum_{r=1}^n X_{i,r}Y_{r,j}$, we can compute substochastic matrices $C^-,C^+$ using $O(d)$-depth, such that
\begin{align*}\left\|C^--\frac1{2^{t-1}}\sum_{i=1}^{2^{t-1}-1}\widetilde D_i\right\|_1&\le\frac1{2^{19d+5}},\\\left\|C^+-\frac1{2^{t-1}}\sum_{i=1}^{2^{t-1}}\widetilde D_i'\right\|_1&\le\frac1{2^{19d+5}}.\end{align*}
Then $2^{t-1}(C^+-C^-)$ is a good approximation of $A^{2^k}$ since
\begin{align*}\|2^{t-1}(C^+-C^-)-A^{2^k}\|_1&\le2^{t-1}\left\|C^--\frac1{2^{t-1}}\sum_{i=1}^{2^{t-1}-1}\widetilde D_i\right\|_1+2^{t-1}\left\|C^+-\frac1{2^{t-1}}\sum_{i=1}^{2^{t-1}}\widetilde D_i'\right\|_1\\&\quad+\sum_{i=1}^{2^{t-1}-1}\|D_i-\widetilde D_i\|_1+\sum_{i=1}^{2^{t-1}}\|D_i'-\widetilde D_i'\|_1\\&\quad+\left\|-\sum_{i=1}^{2^{t-1}-1}D_i+\sum_{i=1}^{2^{t-1}}D_i'-A^{2^k}\right\|_1\\&\le\frac{2^{t-1}}{2^{19d+5}}+\frac{2^{t-1}}{2^{19d+5}}+\frac{2^{t-1}}{2^{19d+5}}+\frac{2^{t-1}}{2^{19d+5}}+\|C-A^{2^k}\|_1\\&\le\frac1{2^{9d+4}}+\left(\sum_{i=1}^{t-1}2^{i-1}\varepsilon_{k-i}^2+2^t\varepsilon_{k-t}\right)\\&\le\frac1{2^{9d+4}}+\frac1{2^{d+2}}.\end{align*}
Here the last step is from the statement of Theorem 19.
Finally we compute a substochastic matrix $C'$ which is a good approximation of $A^{2^k}$ and $2^{t-1}(C^+-C^-)$. Here we need to be careful that $C$ and $2^{t-1}(C^+-C^-)$ are not necessarily non-negative or substochastic (but $A^{2^k}$ is guaranteed substochastic). Let
\begin{align*}C''_{i,j}&:=\max\{2^{t-1}(C_{i,j}^+-C_{i,j}^-),0\},\\C'_{i,j}&:=\frac1{n^{100}}\left[C''_{i,j}\left(1-\frac1{2^{d+1}}\right)\cdot n^{100}\right].\end{align*}
We can compute $C'$ given $C^+,C^-$ by hardwiring the map $(C_{i,j}^+,C_{i,j}^-)\mapsto C_{i,j}'$, which is L-uniform. Obviously $C'$ is non-negative. Note that $C''$ is entrywise closer to $A^{2^k}$ than $2^{t-1}(C^+-C^-)$ and hence
\[\|C''-A^{2^k}\|_1\le\|2^{t-1}(C^+-C^-)-A^{2^k}\|_1\le\frac1{2^{9d+4}}+\frac1{2^{d+2}}.\]
Therefore $C'$ is substochastic since
\[\|C'\|_1\le\left(1-\frac1{2^{d+1}}\right)\|C''\|_1\le\left(1-\frac1{2^{d+1}}\right)\left(1+\frac1{2^{9d+4}}+\frac1{2^{d+2}}\right)\le1.\]
Also note that
\begin{align*}\|C'-A^{2^k}\|_1&\le\|C'-C''\|_1+\|C''-A^{2^k}\|_1\\&\le\frac1{n^{99}}+\frac1{2^{d+1}}\|C''\|_1+\|C''-A^{2^k}\|_1\\&\le\frac1{n^{99}}+\frac1{2^{d+1}}\left(1+\frac1{2^{9d+4}}+\frac1{2^{d+2}}\right)+\frac1{2^{9d+4}}+\frac1{2^{d+2}}\\&\le\frac1{2^d}.\end{align*}
To summarize, we can output a valid $C'$ in $O(d)$-depth. And the circuit is $\operatorname{poly}(n)$-size and L-uniform.
\end{proof}
\section*{5 The Complete Algorithm}
\begin{proof}[Proof of Theorem 20]
Only consider $\log n\ge10$. For $k,t\in\mathbb N$ such that $k\le\log n$ and $1\le t\le3\log n-2k$, we wish to compute a substochastic matrix $M(k,t)$, which is an approximation of $A^{2^k}$, such that $\|M(k,t)-A^{2^k}\|_1\le1/2^t$. Then $M:=M(\log n,\log n)$ is the desired matrix.
For $k=0$, we can trivially let $M(0,t):=A$. Now we show how to recursively compute $M(k_0,t_0)$ for $k_0=1,2,\ldots,\log n$.
In Theorem 18, take the same $n,A$ and take $k:=k_0$, take $B_{k-i}:=M(k-i,[t_0/2]+2i)$ for $1\le i\le k$. Then we can take $\varepsilon_{k-i}:=1/2^{[t_0/2]+2i}$ for $1\le i\le k-1$ and $\varepsilon_0=0$. Now we will invoke Theorem 18,19 by choosing parameter $t$ properly according to the following two cases.
\paragraph{Case 1. $k\le2t_0+2$.}
Take the parameter $t$ in Theorem 18 to be $t:=k$. Then
\[\sum_{i=1}^{k-1}2^{i-1}\varepsilon_{k-i}^2+2^k\varepsilon_0=\sum_{i=1}^{k-1}\frac1{2^{2[t_0/2]+3i+1}}\le\frac1{2^{t_0+2}}.\]
In Theorem 19 take $d:=t_0$. It is easy to verify that $\log n\ge k\ge t$ and $4\log n\ge d\ge t/10$ hold when we invoke Theorem 18,19. Given $B_{k-1},\ldots,B_0$, use $\mathcal I_{n,k_0,k_0,t_0}$ (defined in Theorem 19) we can compute a substochastic matrix $C'$ such that $\|C'-A^{2^k}\|_1\le1/2^{t_0}$.
\paragraph{Case 2. $k\ge2t_0+3$.}
Take $t:=2t_0+3$ in Theorem 18. Then
\[\sum_{i=1}^{2t_0+2}2^{i-1}\varepsilon_{k-i}^2+2^{2t_0+3}\varepsilon_{k-2t_0-3}\le\sum_{i=1}^{2t_0+2}\frac1{2^{2[t_0/2]+3i+1}}+\frac1{2^{[t_0/2]+2t_0+3}}\le\frac1{2^{t_0+2}}.\]
In Theorem 19 take $d:=t_0$. Given $B_{k-1},\ldots,B_0$, use $\mathcal I_{n,k_0,2t_0+3,t_0}$ we can compute a substochastic matrix $C'$ such that $\|C'-A^{2^k}\|_1\le1/2^{t_0}$.
To summarize, take $M(k_0,t_0):=C'$, we can compute $M(k_0,t_0)$ given $M(k_0-i,[t_0/2]+2i)$ for $1\le i\le k_0$, via a $\operatorname{poly}(n)$-size $O(t_0)$-depth L-uniform AC circuit.
Let $\gamma>0$ be a concrete constant such that we can compute $M(k_0,t_0)$ given $M(k_0-i,[t_0/2]+2i)$'s via a $\operatorname{poly}(n)$-size $\gamma t_0$-depth L-uniform AC circuit. Note that if $M(k_0-i,[t_0/2]+2i)$ can be computed in $2\gamma(2(k_0-i)+([t_0/2]+2i))$-depth for $1\le i\le k_0$, then $M(k_0,t_0)$ can be computed in
\[\gamma t_0+\max_{1\le i\le k_0}\{2\gamma(2(k_0-i)+([t_0/2]+2i))\}\le2\gamma(2k_0+t_0)\]
-depth. Also note that $M(0,t_0)$'s are just the inputs, so by induction we know $M(k_0,t_0)$ can be computed in $2\gamma(2k_0+t_0)$-depth. Specially, $M(\log n,\log n)$ (which is the desired output) can be computed in $6\gamma\log n\le O(\log n)$-depth. Also note that we use ``compute $M(k_0,t_0)$ given $M(k_0-i,[t_0/2]+2i)$'' $O((\log n)^2)$ many times, so the total circuit size for computing $M(\log n,\log n)$ is still $\operatorname{poly}(n)$.
\end{proof}

\begin{thebibliography}{99}
\bibitem{source1} AmirMahdi Ahmadinejad, Jonathan A. Kelner, Jack Murtagh, John Peebles, Aaron Sidford, and Salil P. Vadhan. High-precision estimation of random walks in small space. In Sandy Irani, editor, 61st IEEE Annual Symposium on Foundations of Computer Science, FOCS 2020, Durham, NC, USA, November 16-19, 2020, pages 1295–1306. IEEE, 2020. doi:10.1109/ FOCS46700.2020.00123.
\bibitem{source2} Miklós Ajtai. Approximate counting with uniform constant-depth circuits. In Jin-Yi Cai, editor, Advances In Computational Complexity Theory, Proceedings of a DIMACS Workshop, New Jersey, USA, December 3-7, 1990, volume 13 of DIMACS Series in Discrete Mathematics and Theoretical Computer Science, pages 1–20. DIMACS/AMS, 1990. doi:10.1090/DIMACS/ 013/01.
\bibitem{source3} Miklós Ajtai and Michael Ben-Or. A theorem on probabilistic constant depth computations. In Richard A. DeMillo, editor, Proceedings of the 16th Annual ACM Symposium on Theory of Computing, April 30 - May 2, 1984, Washington, DC, USA, pages 471–474. ACM, 1984. doi:10.1145/800057.808715.
\bibitem{source4} Roy Armoni. On the derandomization of space-bounded computations. In International Workshop on Randomization and Approximation Techniques in Computer Science, pages 47–59. Springer, 1998.
\bibitem{source5} László Babai, Noam Nisant, and Márió Szegedy. Multiparty protocols, pseudorandom gen- erators for logspace, and time-space trade-offs. Journal of Computer and System Sciences, 45(2):204–232, 1992.
\bibitem{source6} Mark Braverman, Gil Cohen, and Sumegha Garg. Pseudorandom pseudo-distributions with near-optimal error for read-once branching programs. SIAM Journal on Computing, 49(5):STOC18–242, 2019.
\bibitem{source7} Eshan Chattopadhyay and Jyun-Jie Liao. Optimal error pseudodistributions for read-once branching programs. In 35th Computational Complexity Conference, 2020.
\bibitem{source8} Lijie Chen, William M. Hoza, Xin Lyu, Avishay Tal, and Hongxun Wu. Weighted pseu- dorandom generators via inverse analysis of random walks and shortcutting. In 64th IEEE Annual Symposium on Foundations of Computer Science, FOCS 2023, Santa Cruz, CA, USA, November 6-9, 2023, pages 1224–1239. IEEE, 2023. doi:10.1109/FOCS57990.2023.00072.
\bibitem{source9} Kuan Cheng and William M. Hoza. Hitting sets give two-sided derandomization of small space. In Shubhangi Saraf, editor, 35th Computational Complexity Conference, CCC 2020, July 28-31, 2020, Saarbrücken, Germany (Virtual Conference), volume 169 of LIPIcs, pages 10:1–10:25. Schloss Dagstuhl – Leibniz-Zentrum für Informatik, 2020. doi:10.4230/LIPICS.CCC.2020.10.
\bibitem{source10} Gil Cohen, Dean Doron, Oren Renard, Ori Sberlo, and Amnon Ta-Shma. Error reduction for weighted prgs against read once branching programs. In Valentine Kabanets, editor, 36th Computational Complexity Conference, CCC 2021, July 20-23, 2021, Toronto, Ontario, Canada (Virtual Conference), volume 200 of LIPIcs, pages 22:1–22:17. Schloss Dagstuhl – Leibniz-Zentrum für Informatik, 2021. doi:10.4230/LIPICS.CCC.2021.22.
\bibitem{source11} Gil Cohen, Dean Doron, Oren Renard, Ori Sberlo, and Amnon Ta-Shma. Error reduction for weighted prgs against read once branching programs. In Proceedings of the 36th Computational Complexity Conference, page 1, 2021.
\bibitem{source12} Gil Cohen, Dean Doron, Ori Sberlo, and Amnon Ta-Shma. Approximating iterated multiplic- ation of stochastic matrices in small space. In Barna Saha and Rocco A. Servedio, editors, Proceedings of the 55th Annual ACM Symposium on Theory of Computing, STOC 2023, Or- lando, FL, USA, June 20-23, 2023, pages 35–45. ACM, 2023. doi:10.1145/3564246.3585181.
\bibitem{source13} Joshua Cook. Size bounds on low depth circuits for promise majority. In Nitin Saxena and Sunil Simon, editors, 40th IARCS Annual Conference on Foundations of Software Technology and Theoretical Computer Science, FSTTCS 2020, December 14-18, 2020, BITS Pilani, K K Birla Goa Campus, Goa, India (Virtual Conference), volume 182 of LIPIcs, pages 19:1–19:14. Schloss Dagstuhl – Leibniz-Zentrum für Informatik, 2020. doi:10.4230/LIPICS.FSTTCS.2020.19.
\bibitem{source14} Dean Doron, Edward Pyne, and Roei Tell. Opening up the distinguisher: A hardness to randomness approach for BPL = L that uses properties of BPL. Electron. Colloquium Comput. Complex., TR23-208, 2023. arXiv:TR23-208.
\bibitem{source15} Dean Doron and Roei Tell. Derandomization with minimal memory footprint. In Amnon Ta-Shma, editor, 38th Computational Complexity Conference, CCC 2023, July 17-20, 2023, Warwick, UK, volume 264 of LIPIcs, pages 11:1–11:15. Schloss Dagstuhl – Leibniz-Zentrum für Informatik, 2023. doi:10.4230/LIPICS.CCC.2023.11.
\bibitem{source16} Pooya Hatami and William Hoza. Theory of unconditional pseudorandom generators. Electron. Colloquium Comput. Complex., TR23-019, 2023. arXiv:TR23-019.
\bibitem{source17} William M. Hoza. Better pseudodistributions and derandomization for space-bounded com- putation. In Mary Wootters and Laura Sanità, editors, Approximation, Randomization, and Combinatorial Optimization. Algorithms and Techniques, APPROX/RANDOM 2021, August 16-18, 2021, University of Washington, Seattle, Washington, USA (Virtual Conference), volume
\bibitem{source207} of LIPIcs, pages 28:1–28:23. Schloss Dagstuhl – Leibniz-Zentrum für Informatik, 2021. doi:10.4230/LIPICS.APPROX/RANDOM.2021.28.
\bibitem{source18} William M. Hoza. Recent progress on derandomizing space-bounded computation. Bull. EATCS, 138, 2022. URL: http://eatcs.org/beatcs/index.php/beatcs/artic$\ell$e/view/728.
\bibitem{source19} William M Hoza and David Zuckerman. Simple optimal hitting sets for small-success rl. SIAM Journal on Computing, 49(4):811–820, 2020.
\bibitem{source20} Russell Impagliazzo, Noam Nisan, and Avi Wigderson. Pseudorandomness for network algorithms. In Proceedings of the twenty-sixth annual ACM symposium on Theory of computing, pages 356–364, 1994.
\bibitem{source21} Adam R. Klivans and Dieter van Melkebeek. Graph nonisomorphism has subexponential size proofs unless the polynomial-time hierarchy collapses. SIAM J. Comput., 31(5):1501–1526, 2002. doi:10.1137/S0097539700389652.
\bibitem{source22} Noam Nisan. Pseudorandom generators for space-bounded computation. Comb., 12(4):449–461, 1992. doi:10.1007/BF01305237.
\bibitem{source23} Noam Nisan. RL $\subseteq$ SC. Comput. Complex., 4:1–11, 1994. doi:10.1007/BF01205052.
\bibitem{source24} Noam Nisan and David Zuckerman. Randomness is linear in space. Journal of Computer and System Sciences, 52(1):43–52, 1996.
\bibitem{source25} Aaron (Louie) Putterman and Edward Pyne. Near-optimal derandomization of medium-width branching programs. In Barna Saha and Rocco A. Servedio, editors, Proceedings of the 55th Annual ACM Symposium on Theory of Computing, STOC 2023, Orlando, FL, USA, June 20-23, 2023, pages 23–34. ACM, 2023. doi:10.1145/3564246.3585108.
\bibitem{source26} Edward Pyne, Ran Raz, and Wei Zhan. Certified hardness vs. randomness for log-space. In 64th IEEE Annual Symposium on Foundations of Computer Science, FOCS 2023, Santa Cruz, CA, USA, November 6-9, 2023, pages 989–1007. IEEE, 2023. doi:10.1109/FOCS57990.2023.00061.
\bibitem{source27} Edward Pyne and Salil Vadhan. Pseudodistributions that beat all pseudorandom generators. In 36th Computational Complexity Conference (CCC 2021). Schloss Dagstuhl – Leibniz-Zentrum für Informatik, 2021.
\bibitem{source28} Edward Pyne and Salil P. Vadhan. Pseudodistributions that beat all pseudorandom generators (extended abstract). In Valentine Kabanets, editor, 36th Computational Complexity Conference, CCC 2021, July 20-23, 2021, Toronto, Ontario, Canada (Virtual Conference), volume 200 of LIPIcs, pages 33:1–33:15. Schloss Dagstuhl – Leibniz-Zentrum für Informatik, 2021. doi: 10.4230/LIPICS.CCC.2021.33.
\bibitem{source29} Michael E. Saks and Shiyu Zhou. BPH SPACE(S) $\subseteq$ DSPACE(S 3/2 ). J. Comput. Syst. Sci., 58(2):376–403, 1999. doi:10.1006/JCSS.1998.1616.
\bibitem{source30} Emanuele Viola. On approximate majority and probabilistic time. In 22nd Annual IEEE Conference on Computational Complexity (CCC 2007), 13-16 June 2007, San Diego, California, USA, pages 155–168. IEEE Computer Society, 2007. doi:10.1109/CCC.2007.16.
\bibitem{source31} Emanuele Viola. Randomness buys depth for approximate counting. In Rafail Ostrovsky, editor, IEEE 52nd Annual Symposium on Foundations of Computer Science, FOCS 2011, Palm Springs, CA, USA, October 22-25, 2011, pages 230–239. IEEE Computer Society, 2011. doi:10.1109/FOCS.2011.19.
\end{thebibliography}
\end{document}
